% Initial common-annihilator calculation proposed by /root and independently
% reconstructed here. GPT-6 Astra Pro; effort not exposed;
% /root/replacement; 22 November 2026.
\section{HD-F1-R11: the balanced kernel determines every retract character and excludes torsion endpoints}
\label{sec:followup-rm-character-rigidity}

Use the genuine real-quadratic fourfold, compatible polarization,
and full character of R07:
\[
 v=(1-qA_1)(1-qA_2),\quad A_i=\theta_i^2/2,\quad
 \operatorname{ch}(E)=N e^{-B}v,\quad N\ne0,\quad q>0.
 \tag{F1R11.1}
\]
Every actual weakly $G$-semiregular lift for a coherent faithful
finite point action has the following stronger properties:
\[
 \begin{split}
 &\sigma_E^0\bigl(\operatorname{End}_X(E)\bigr)
                   =\mathbf C\operatorname{ch}(E);\\
 &\operatorname{ch}(Q)\in\mathbf C\operatorname{ch}(E)
       \quad\text{for every actual retract }Q\text{ of }E.
 \end{split}
 \tag{F1R11.2}
\]
These concern the full ordinary endomorphism algebra and every
ordinary retract, including ones not preserved by $G$. They need
neither the minimal condition $e_1^G=8$ nor negative-Ext vanishing.
For rational characters, the proportionality scalar for a retract
is rational, because it is the ratio of the two ranks.

There is also an unconditional obstruction for actual coupled
objects. A nonzero torsion sheaf cannot be an extreme standard
cohomology sheaf of such an $E$ separated from the remaining
cohomology by a zero adjacent degree. More generally, every
extreme cohomology sheaf across such a gap must have full
character proportional to (F1R11.1), with the cohomological
shift sign included.

In particular no actual two-stage object
\[
 A\longrightarrow E\longrightarrow T[-m]\longrightarrow A[1],
 \qquad A,T\text{ coherent},\quad T\ne0\text{ torsion},\quad m\ge2,
 \tag{F1R11.3}
\]
whose whole character is (F1R11.1), has a weakly semiregular
lift for any of these actions. The coupling class in
$\operatorname{Ext}^{m+1}(T,A)$ is arbitrary. This excludes
the entire torsion-endpoint amplitude $[0,2]$ family of
DEG\_L S-27/S-29, in every positive virtual rank for which
the balanced whole-character equation is posed. No simplicity,
principal polarization, smoothness of a support, or vanishing
of its special cohomology is needed.

\paragraph{1. The common annihilator in the complete HKR space.}
The normalized degree-two kernel is
$K=K_1\oplus K_2$ from (F1R7.6). We now compute its common
annihilator, rather than only the ranks of its action on $v$.
In one eigenspace block choose holomorphic generators $x_1,x_2$
and antiholomorphic generators $y_1,y_2$ such that
$A_i=x_1x_2y_1y_2$. Such a rescaling is possible because $A_i$
is a nonzero product of the two volume forms. Put
$t_a=\iota_{x_a^*}$ and $w_b=y_b\wedge$. With the composition
convention $t_1t_2(A_i)=-y_1y_2$, a basis of $K_i$ acts as
\[
 t_aw_b\quad(1\le a,b\le2),\qquad
                  t_1t_2-q\,w_1w_2.
 \tag{F1R11.4}
\]
The four mixed operators annihilate exactly the seven-dimensional
subspace with basis
\[
 1,\quad y_1,\quad y_2,\quad y_1y_2,\quad
 x_1y_1y_2,\quad x_2y_1y_2,\quad x_1x_2y_1y_2.
 \tag{F1R11.5}
\]
For clarity, write a general form by holomorphic and
antiholomorphic degrees. Terms of holomorphic degree zero
are unrestricted. Terms of positive holomorphic degree and
antiholomorphic degree two are also killed. For a term of
positive holomorphic degree and antiholomorphic degree zero
or one, contraction by each $x_a^*$ followed by wedge by
both $y_b$ forces every coefficient to vanish. The operator
bidegrees separate these equations, proving precisely
(F1R11.5).

The last operator in (F1R11.4) sends $c+dA_i$ to
$-(qc+d)y_1y_2$ and kills the other five displayed basis
vectors. Its common kernel $M_i$ therefore has the basis
\[
 1-qA_i,\quad y_1,\quad y_2,\quad y_1y_2,\quad
 x_1y_1y_2,\quad x_2y_1y_2.
 \tag{F1R11.6}
\]
Its Hochschild homological degrees (holomorphic minus
antiholomorphic degree) are respectively zero, minus one
four times, and minus two once.

The full HKR space is the exterior tensor product of the
two blocks. Operators in $K_1$ and $K_2$ act on the respective
factors, so their common kernel is exactly $M_1\otimes M_2$.
All its block degrees are nonpositive. Its total degree-zero
part is consequently the single line generated by the
tensor product of the two degree-zero vectors:
\[
 \{z\in HH_0(X):a\cdot z=0\text{ for all }a\in K\}
                    =\mathbf C v.
 \tag{F1R11.7}
\]
This checks the whole seventy-dimensional $HH_0(X)$, including
off-block Hodge components; it does not assume that a
character is a polynomial in $\theta_1,\theta_2$.
The shear (F1R7.5) gives the unnormalized identity
\[
 \{z\in HH_0(X):a\cdot z=0\text{ for all }a\in K_E\}
                    =\mathbf C\operatorname{ch}(E),
 \quad K_E=\ker c^2_{\operatorname{ch}(E)}.
 \tag{F1R11.8}
\]

\paragraph{2. Weak semiregularity and the full ordinary trace.}
Every $a\in K_E$ has invariant characteristic evaluation and
zero full trace. The actual identity-sector comparison in R05
identifies this kernel with the equivariant trace kernel.
Weak semiregularity therefore gives
\[
                  \operatorname{ev}_E(a)=0
                  \quad\text{in ordinary }\operatorname{Ext}^2_X(E,E).
 \tag{F1R11.9}
\]
The ordinary characteristic algebra map consequently factors
through the transformed exterior quotient $Q$ in (F1R7.7).
Its degree-four image has dimension at most one. The
rank-one map $c^4_{\operatorname{ch}(E)}
=\sigma_E^4\operatorname{ev}_E^4$ makes that dimension
exactly one. Ordinary CY4 Serre duality and the
Buchweitz--Flenner characteristic/trace transpose now give
rank one for the ordinary map $\sigma_E^0$. Its value on
the identity is the nonzero whole character. This proves
the first equality in (F1R11.2) without any condition on
$e_1^G$ or on the dimensions of ordinary self-Ext groups.

For an actual retract with inclusion $i$ and projection $p$,
the perfect trace identity is
$\sigma_E^0(ip)=\operatorname{ch}(Q)$. It proves the
second assertion of (F1R11.2). There is also an independent
proof using (F1R11.8): naturality gives
$p[2]\operatorname{ev}_E(a)i=\operatorname{ev}_Q(a)=0$
for every $a\in K_E$, hence
$c^2_{\operatorname{ch}(Q)}(K_E)=0$. The complete
annihilator calculation forces the same proportionality.

\paragraph{3. The invariant endomorphism conclusion and its precise remaining scope.}
Restrict the rank-one ordinary trace to
$A_G=\operatorname{End}_G(E)$. There is a unique functional
$\lambda:A_G\to\mathbf C$ with
\[
 \sigma_E^0(u)=\lambda(u)\operatorname{ch}(E),\qquad
                         \lambda(1)=1.
 \tag{F1R11.10}
\]
For every weakly semiregular lift, its invariant degree-two
evaluation is surjective and $e_2^G=18$, by the same actual
Serre transpose as in R07. Products of two such evaluations
belong to the one-dimensional degree-four characteristic
image. Pairing with $u$ and using (F1R11.10), invariant
Serre duality gives
\[
 ux=xu=\lambda(u)x
             \quad(x\in\operatorname{Ext}^2_G(E,E)).
 \tag{F1R11.11}
\]
For a nonzero such $x$, associativity shows
$\lambda(uv)=\lambda(u)\lambda(v)$. Thus $\lambda$ is a
unital algebra character, and $\ker\lambda$ is a
codimension-one two-sided ideal annihilating invariant
degree two on both sides.

This conclusion concerns $A_G$. Multiplicativity is not
asserted for the analogous functional on the full ordinary
endomorphism algebra. Moreover annihilation of invariant
degrees one and three and locality of $A_G$ use precisely
the additional assumptions $e_1^G=8$ and invariant
negative-Ext vanishing in R07. With these retained, R07
proves its two local three-dimensional possibilities.
No claim about those two algebras' geometric existence
has been added here.

\paragraph{4. Actual truncation maps exclude every torsion endpoint across a gap.}
Let $Q$ be either shifted extreme cohomology sheaf of $E$
in (F1R10.4) or (F1R10.6), with the zero adjacent
cohomology degree as there. The exact injections
(F1R10.5) and (F1R10.7) and naturality show that
(F1R11.9) forces
\[
              \operatorname{ev}_Q(K_E)=0.
 \tag{F1R11.12}
\]
The trace--Atiyah identity and (F1R11.8) then give
$\operatorname{ch}(Q)\in\mathbf C\operatorname{ch}(E)$.
This proves the stated character constraint on every
extreme cohomology sheaf across a gap, without a
pure-character hypothesis on that sheaf.

A nonzero coherent torsion sheaf $T$ has rank zero and
nonzero full character. To check the latter explicitly,
let $d$ be the maximum dimension of its support. Its
first nonzero character component is the effective cycle
\[
 \operatorname{ch}_{4-d}(T)
   =\sum_{\dim Z=d}\operatorname{length}_{\mathcal O_{X,\eta_Z}}
                                  (T_{\eta_Z})[Z].
 \tag{F1R11.13}
\]
Indeed, in the Grothendieck group modulo sheaves of smaller
support dimension, the generic-length filtration identifies
$[T]$ with the displayed sum of $[\mathcal O_Z]$. On the
smooth dense part of each $Z$, the regular-immersion GRR
formula is
$\operatorname{ch}(j_*\mathcal O_Z)=j_*\operatorname{td}(N_{Z/X})^{-1}$,
whose leading term is $[Z]$. The omitted smaller-dimensional
support cannot contribute in codimension $4-d$.
Every coefficient is positive. Pairing this component
with an ample class to the power $d$ gives a strictly
positive number, so it is nonzero in cohomology.
Since $\operatorname{ch}(E)$ has nonzero rank, no
nonzero rank-zero character can be proportional to it.
The shift of $T$ only multiplies its character by a
sign. Thus (F1R11.12) is impossible for a torsion endpoint.

Equivalently this proof constructs a nonzero kernel
element: choose one of the ten normalized operators in
(F1R11.4), transported back by the explicit $B$-shear,
whose action on $\operatorname{ch}(T)$ is nonzero.
Such an operator exists by (F1R11.8) and (F1R11.13).
Its evaluation on $T$ is nonzero; the actual injection
through $i$ or $p$ makes its evaluation on $E$ nonzero.
The whole character annihilates it, and every coherent
point lift fixes it. This is the nonzero kernel of the
original weak trace map.

For (F1R11.3) the injection is supplied simply by
$\operatorname{Ext}^{1-m}_X(A,T)=0$, as computed in
(F1R10.11). In the specific surface-support construction
S-29, $T=j_*M$ has the explicit leading cycle
$\operatorname{ch}_2(T)=[S]\ne0$. Thus every one of its
actual nonzero coupling classes is covered by the
preceding proof, independent of the special-cohomology
conditions listed in S-29.

\paragraph{Framework and checking scope.}
The exterior kernel here is the actual HH action entering
the published Buchweitz--Flenner diagram; all ten operators
and their complete common annihilator are written above.
The Grothendieck-group map of the actual triangle is
$[E]=[A]+(-1)^m[T]$, and (F1R11.13) is its leading
cycle realization by GRR. Adams operations and the
Poincar\'e biextension/finite-action coherences are the
same precise morphisms (F1R9.9)--(F1R9.12). The original
common-annihilator idea was proposed by the root agent
and independently checked by the full block computation
in paragraph 1; the ordinary trace proof in paragraph 2
also follows independently from the previously proved
exterior quotient and the ordinary Serre transpose.
